\chapter{Results}
\label{ch:results}

This chapter presents the empirical results of the statistical benchmarking suite, evaluating MalthusJAX across the three core hypotheses formulated in Chapter \ref{ch:methodology}: Mathematical Parity (H1), Structural Ablation (H2), and Numerical Representation (H3). The experiments systematically isolate the computational overhead and scaling behaviors of the framework as problem dimensionality $D$ scales from 10 to 1,000.

% ==============================================================================
% 5.1 MATHEMATICAL PARITY (H1)
% ==============================================================================
\section{Mathematical Parity (H1)}
\label{sec:results_parity}

The first requirement of MalthusJAX is to serve as a mathematically identical drop-in replacement for the native Python framework, EvoSAX. To test this, the \texttt{malthusjax\_wrapper} pipeline was executed against the \texttt{evosax\_baseline} across the BBOB testbed (Sphere, Rosenbrock, Rastrigin).

\subsection{Non-Parametric Location Shifts}
Because execution times and fitness landscapes are heavily skewed at varying dimensionalities, we utilize the non-parametric Wilcoxon Signed-Rank test to evaluate location shifts between the paired solvers. Table \ref{tab:wilcoxon_parity} presents the summary statistics and computed $p$-values.

\begin{table}[h]
    \centering
    \caption{Wilcoxon Parity Results for Baseline Equivalence.}
    \label{tab:wilcoxon_parity}
    \resizebox{\textwidth}{!}{
        % INJECT PARITY TABLE ARTIFACT
        \input{results/h1_parity_lhs_smoke/analysis/ols_regression_table.tex}
    }
\end{table}

The computed Cohen's $d_z$ effect sizes are infinitesimally small (e.g., $-0.04$ for the Sphere execution time), indicating that the structural behavior is statistically indistinguishable.

\subsection{Log-Log Scaling Laws}
To guarantee that the JAX JIT compiler overhead does not introduce a hidden asymptotic complexity penalty, we constructed a Log-Log Interaction Regression model:
\begin{equation}
    \ln(\text{Metric}) = \beta_0 + \beta_1 I_{mjx} + \beta_2 \ln(D) + \beta_3 (I_{mjx} \times \ln(D))
\end{equation}

Table \ref{tab:ols_parity} proves that the interaction coefficient ($\beta_3$) is universally non-significant across all execution time scaling regressions (e.g., $p=0.95$ for Rastrigin, $p=0.81$ for Sphere). Consequently, we reject the presence of any asymptotic deviation; MalthusJAX scales identically to EvoSAX.

\begin{table}[h]
    \centering
    \caption{Log-Log OLS Scaling Regression Coefficients.}
    \label{tab:ols_parity}
    \resizebox{\textwidth}{!}{
        % INJECT OLS TABLE ARTIFACT
        \input{results/h1_parity_lhs_smoke/analysis/ols_regression_table.tex}
    }
\end{table}

\begin{figure}[h]
    \centering
    \includegraphics[width=0.8\textwidth]{results/h1_parity_lhs_smoke/analysis/malthusjax_wrapper_vs_evosax_baseline_sphere_execution_time_scaling.png}
    \caption{Log-Log scaling of Execution Time vs Dimensionality ($D$) on the Sphere function.}
    \label{fig:parity_sphere_scaling}
\end{figure}

% ==============================================================================
% 5.2 STRUCTURAL ABLATION (H2)
% ==============================================================================
\section{Structural Ablation (H2)}
\label{sec:results_ablation}

Having established baseline parity, we now dissect the specific computational overhead introduced by individual operator compilation. Table \ref{tab:ols_ablation} highlights the differential cost of substituting the wrapped EvoSAX routines with pure JAX-native MalthusJAX operators (Sparse Gaussian Mutation, Uniform Real Crossover, Tournament Selection, Elite Pool).

\begin{table}[h]
    \centering
    \caption{Ablation OLS Regression (Native vs Wrapper).}
    \label{tab:ols_ablation}
    \resizebox{\textwidth}{!}{
        % WILL INJECT ABLATION TABLE ARTIFACT ONCE RUN IS COMPLETE
        \input{results/h2_ablation_smoke/analysis/ols_regression_table.tex}
    }
\end{table}

\begin{figure}[h]
    \centering
    \includegraphics[width=0.8\textwidth]{results/h2_ablation_smoke/analysis/mjx_ablate_mutation_vs_mjx_baseline_sphere_execution_time_scaling.png}
    \caption{Scaling impact of native Sparse Gaussian Mutation.}
    \label{fig:ablation_mutation}
\end{figure}

% ==============================================================================
% 5.3 NUMERICAL REPRESENTATION (H3)
% ==============================================================================
\section{Numerical Representation (H3)}
\label{sec:results_representation}

Finally, we evaluate the system's ability to maintain solver convergence fidelity while aggressively down-casting tensor precision. The execution times and best fitness deviations for \texttt{float32}, \texttt{bfloat16}, and \texttt{float16} are benchmarked.

\begin{table}[h]
    \centering
    \caption{Representation Parity (Mixed Precision).}
    \label{tab:wilcoxon_representation}
    \resizebox{\textwidth}{!}{
        % WILL INJECT REPRESENTATION TABLE ARTIFACT ONCE RUN IS COMPLETE
        \input{results/h3_representation_smoke/analysis/ols_regression_table.tex}
    }
\end{table}

\begin{figure}[h]
    \centering
    \includegraphics[width=0.8\textwidth]{results/h3_representation_smoke/analysis/mjx_bf16_vs_mjx_f32_sphere_boxplots.png}
    \caption{Location shifts comparing f32 and bf16 on Rosenbrock.}
    \label{fig:representation_bf16_box}
\end{figure}
